Estimating correlation matrices between financial assets from historical series alone suffers from two structural limitations: the statistical noise that degrades empirical matrices when the number of assets is of the same order of magnitude as the length of the series, and the fact that the past offers a single trajectory among all those that were possible. This thesis, carried out in collaboration with Intesa Sanpaolo, addresses both limitations by testing the effectiveness of Variational Autoencoders (VAE) as stochastic generative models for the simulation of realistic market scenarios.

The work starts from the daily returns of 362 S\&P 500 constituents continuously quoted over the twenty-year period 2000-2020. From them, using a rolling window of 724 days with a stride of 10, 461 empirical correlation matrices are extracted, reduced to 412 by a validity filter based on the Perron-Frobenius property. Since a matrix produced by a generative model is not automatically positive semidefinite --- a necessary condition for any valid correlation matrix --- the input to the models is not the raw matrix but the lower triangular factor of its Cholesky decomposition: any vector produced by the decoder therefore recomposes into a positive semidefinite matrix by construction.

The methodology follows an incremental path across four architectures compared at equal latent dimension: Principal Component Analysis as a linear benchmark, Linear Autoencoder, Deep Autoencoder and Variational Autoencoder. The evaluation combines spatial error metrics (MSE, MAE, Frobenius norm) with topological metrics built on the Minimum Spanning Tree. On pure historical reconstruction the Deep Autoencoder achieves the best results at every latent dimension considered; the VAE, trained with a reconstruction loss based on the Jeffreys Divergence between the distributions associated with the matrices, sits in an intermediate position: at twenty latent dimensions it outperforms both PCA and the Linear Autoencoder, while remaining behind the deterministic Autoencoder. Variational regularisation therefore does not substantially compromise reconstruction fidelity.

It is on the generative side that the VAE shows its distinctive value: the continuity and density of the latent space induced by the Kullback-Leibler Divergence are not guaranteed in a deterministic autoencoder, whose latent space remains fragmented. The procedure adopted is not a simple sampling from the prior, but a historical block bootstrap of the daily latent increments, which anchors each scenario to the market state observed today and produces a distribution of plausible one-month evolutions of the correlation structure, to be assessed from a risk perspective.

The one thousand generated scenarios were validated against the five stylized facts of financial correlation matrices: a coefficient distribution shifted towards positive values, an eigenvalue spectrum consistent with the Marchenko-Pastur law, the Perron-Frobenius property, hierarchical sector structure and the scale-free topology of the Minimum Spanning Tree. All one thousand scenarios satisfy the five criteria simultaneously, with cosine similarity above 0.98 with respect to today's matrix on the principal modes. The proposed architecture thus emerges as a self-contained and structurally coherent tool for generating realistic synthetic financial ecosystems.
